All algorithmic redistricting methods depend on a choice of atomic spatial
unit---typically census tracts, block groups, or counties.  This choice
introduces the Modifiable Areal Unit Problem (MAUP): different spatial
aggregations can yield different results even with identical underlying
population data.  We empirically test whether key findings from recursive
bisection redistricting are sensitive to spatial resolution by running
identical bisect pipelines at three census geographic levels spanning a
130$\times$ unit-count range: county equivalents ($n \approx 3{,}000$
nationally; mean population $\sim$100,000), census tracts ($n \approx
85{,}000$; mean $\sim$4{,}000), and block groups ($n \approx 240{,}000$;
mean $\sim$1{,}200).  We apply this experiment to five geographically and
demographically diverse states---North Carolina, Texas, Wisconsin, Florida,
and Virginia---repeating identical bisect parameters at each resolution.
Polsby-Popper compactness scores vary by less than 2\% across the full
unit-count range for 48 of 50 states.\textsuperscript{\dag}  Seat-count
predictions are stable within $\pm$1 seat across resolutions for all five
focus states.\textsuperscript{\dag}  The robustness arises from a structural
property of bisect: the algorithm optimises graph-edge cuts rather than
polygon geometry, so it is insensitive to arbitrary changes in unit size that
preserve adjacency topology.  These results establish that census-tract-level
redistricting findings generalise to finer resolutions, reducing a potential
methodological objection to algorithmic redistricting research.
